\section{Discussion}
\label{sec:discussion}

\subsection{Why PSO Did Not Beat Random Search}
\label{subsec:why}

We believe there are several reasons for this outcome. In \genone, the default values appear to lie in a wide and flat region of good settings. They scored higher than most of the proxy evaluations of both searches and were behind the two winners by only 0.019 (PSO) and 0.023 (random search), as reported in Section~\ref{subsec:res-search}. In the proxy check they were second out of eight at both training lengths, and after 40 epochs the four best settings were within 0.026 of each other (Table~\ref{tab:proxy}).

Inside such a region the proxy is able to distinguish poor settings from good ones ($\rho = 0.976$), but it cannot reliably identify the best one. The setting that won on the proxy in \genone, that of random search, actually gave the lowest validation score after 100 epochs (Table~\ref{tab:gen1-test}). A short proxy may favor settings that learn quickly. In addition, taking the maximum over 32 single-run scores tends to select a score that is high partly by chance. Neither of these effects was measured by us.

The budget was small as well. There were 32 evaluations for six dimensions, and PSO used eight of them for its random initial swarm. Bergstra and Bengio~\cite{bergstra2012random} showed that for most datasets only a few hyperparameters really matter, and that random search finds models that are as good as or better than those of grid search for the same budget. If one or two of our six dimensions dominate, then 32 independent samples already cover them quite densely and little is left for a guided search to exploit. The fact that all four search winners have an initial learning rate below the default is consistent with this explanation, although we did not measure the importance of each dimension.

The swarm also converged early. Its mean proxy score became higher than that of random search, but its best score did not (Section~\ref{subsec:res-search}). In \gentwo, the best particle was the last one of the random initial swarm (Fig.~\ref{fig:convergence}). In effect, PSO was offering the best of eight random samples against the best of 32. With only three update steps, a global-best swarm gathers its later evaluations around the best point of round~1. This increases the mean but not the maximum, and only the maximum is used.

It is less clear why tuning helped in \gentwo alone. Our hypothesis is that the new data (real-world photographs, 640~px images and a seventh class) shifted the good region away from the defaults, so that any reasonable search could gain about 0.02 \mapfifty on the mixed and studio test sets. This could not be verified because there is no proxy run of the defaults for \gentwo. As every arm was trained only once, some of the gain may also be training noise. Repeating the proxy check and the full trainings for \gentwo would help to separate the two explanations (Section~\ref{subsec:fw-opt}).

\subsection{Implications for Tuning Claims in Waste AI}
\label{subsec:implications}

According to Gap~3 (Section~\ref{subsec:gaps}), metaheuristic tuning in this area is compared with other metaheuristics or with nothing, and single test values are reported. Our results illustrate what can go wrong in that case. If PSO had been compared with the defaults alone, its significant gain of $+0.020$ \mapfifty in \gentwo would have been attributed to the swarm. Random search, however, gained $+0.017$ with the same budget, and the two tuned arms are indistinguishable (Table~\ref{tab:gen2-test}). The gain is therefore a result of searching as such and not of the swarm. Single values would have been misleading too. In \genone the point estimates place the defaults first and PSO last, and in \gentwo the validation scores place random search ahead of PSO, but neither order is significant on the test set. In the opposite direction, the significant $+0.003$ of PSO over random search on the India set, where scores are about 0.02, shows that a significant difference also has to be large enough to matter in practice.

On this basis, we suggest the following minimum protocol for any claim that a metaheuristic tunes a waste model better: (1) include the defaults and random search with the same number of evaluations, and report the time taken (Section~\ref{subsec:res-cost}); (2) validate any cheap proxy against longer training before the search, using a threshold decided beforehand; (3) select the final setting of each arm on validation data and evaluate the test split once; (4) report paired confidence intervals for the differences between arms on each test domain, and judge them against a difference that is relevant in practice; and (5) repeat the trainings and searches with several seeds. SwarmSort satisfies all of these except the fifth. For equal-budget comparisons such as the one by Sayed et al.~\cite{sayed2024intelligent}, adding a random-search arm would cost no more than one additional competitor.

\subsection{Real-World and Indian Deployment}
\label{subsec:deployment}

For a smart bin or a municipal service, the data are far more important than the optimizer. Every arm loses about half of its studio score on real-world photographs and drops to 0.023--0.027 \mapfifty on the Indian ones (Section~\ref{subsec:res-domain}). The largest effect of tuning, about 0.02, is less than one tenth of the drop from studio to real-world images (about 0.28), and it was not observed on real-world photographs at all. The strongest evidence comes from Table~\ref{tab:g1g2}, where both generations are evaluated on the same images. Real-world training data increased the real-world \mapfifty from 0.038--0.096 in \genone to 0.338--0.353 in \gentwo, with every \gentwo arm significantly above every \genone arm. Tuning, in contrast, made no significant difference on real-world images in either generation.

This improvement came with a significant loss of 0.074 to 0.105 on studio images. Three causes are likely, and our runs cannot separate them. \gentwo retained only 41\% of the studio training images of \genone. The same small YOLOv8n now has to share its capacity with real-world scenes and a seventh class. Finally, \gentwo was trained at 640~px instead of 416~px and its checkpoint was selected on mixed validation data. We conclude that local training data will be of more help than a better search, and this is why Indian data are placed first in our future work (Section~\ref{subsec:fw-data}). Clutter is the second difficulty. The models label all objects correctly in about 44\% of the real-world photographs that contain one object, but in none of the 46 photographs with three or more objects (Fig.~\ref{fig:best}). A bin that photographs one item at a time is thus a realistic goal at present, whereas sorting a mixed pile is not.

Since the merged, mask-derived boxes of DWSD penalize correct per-item detections (Section~\ref{subsec:res-domain}), a test set needs to be labeled in the same manner as the training data before its score can be interpreted as accuracy. Relabeling a part of DWSD item by item would make it possible to separate the labeling style from the model error.

The Indian rules also have a bearing on the classes. Rule~15(h) of the Solid Waste Management Rules, 2016~\cite{moefcc2016solid} prescribes green, white and black bins for bio-degradable, recyclable and other waste. BIODEGRADABLE corresponds to the green bin, and the material classes correspond to the white bin to the extent that their items are recyclable. The black bin corresponds to OTHER, which is the weakest class of every arm (AP@50 of 0.050 to 0.063). A bin built on these models would be least reliable exactly there.

\subsection{Limitations and Threats to Validity}
\label{subsec:limitations}

This study has a number of limitations. Each arm was trained once (seed 42) and each search was run once (seed 42 for PSO and seed 123 for random search). The confidence intervals therefore cover only the sampling of test images, and the small differences between arms could change with another seed. We did not compute $p$-values or correct for multiple comparisons, although \gentwo involves twelve pairwise comparisons per metric (three pairs on four test sets). Per-class AP has no interval, and the single BIODEGRADABLE box in the real-world test set makes the intervals of that set wide and asymmetric.

The proxy check was based on eight settings, was done for \genone only and used a 40-epoch reference instead of a 100-epoch one, so the validity of the proxy in \gentwo is assumed. The proxy evaluated the weights of the last epoch, while the final models used the best validation checkpoint (\texttt{best.pt}). The Arm~C training of \gentwo was resumed after epoch 55 because of a frozen session. It kept the same weights and optimizer state but is not bit-identical to an uninterrupted run.

For DWSD, the class mapping was inferred from the masks with six of the fifteen mask values uncertain, and the boxes are connected regions and not individual items. Only six training hyperparameters were searched, and the budget was equal in evaluations but not in time. We did not run any external baseline, such as the pipeline of Majchrowska et al.~\cite{majchrowska2022deep}, on our data, so our scores cannot be ranked against earlier work. Lastly, the literature review was targeted and not systematic (Section~\ref{subsec:review-method}), and the gaps in Section~\ref{subsec:gaps} hold only for the literature that we reviewed.
